\documentclass[11pt]{article}

\usepackage[T1]{fontenc}
\usepackage[utf8]{inputenc}
\usepackage{lmodern}
\usepackage[margin=1.02in]{geometry}
\usepackage{amsmath,amssymb,amsthm,mathtools}
\usepackage{booktabs}
\usepackage{array}
\usepackage{longtable}
\usepackage{tabularx}
\usepackage{enumitem}
\usepackage{xcolor}
\usepackage{hyperref}

% Lean declarations contain long alphanumeric components.  Permit line
% breaks inside those components when they are typeset with \path.
\makeatletter
\g@addto@macro\UrlBreaks{%
  \do\a\do\b\do\c\do\d\do\e\do\f\do\g\do\h\do\i\do\j\do\k\do\l\do\m%
  \do\n\do\o\do\p\do\q\do\r\do\s\do\t\do\u\do\v\do\w\do\x\do\y\do\z%
  \do\A\do\B\do\C\do\D\do\E\do\F\do\G\do\H\do\I\do\J\do\K\do\L\do\M%
  \do\N\do\O\do\P\do\Q\do\R\do\S\do\T\do\U\do\V\do\W\do\X\do\Y\do\Z}
\makeatother

\hypersetup{
  colorlinks=true,
  linkcolor=blue!55!black,
  citecolor=green!35!black,
  urlcolor=blue!55!black
}
\setlist{itemsep=2pt,topsep=4pt}
\setlength{\parindent}{0pt}
\setlength{\parskip}{5pt}
\setlength{\emergencystretch}{3em}
\renewcommand{\arraystretch}{1.16}

\newtheorem{theorem}{Theorem}[section]
\newtheorem{proposition}[theorem]{Proposition}
\newtheorem{lemma}[theorem]{Lemma}
\newtheorem{corollary}[theorem]{Corollary}
\newtheorem{definition}[theorem]{Definition}
\newtheorem{remark}[theorem]{Remark}
\newtheorem{warning}[theorem]{Warning}

\newcommand{\R}{\mathbb R}
\newcommand{\Q}{\mathbb Q}
\newcommand{\N}{\mathbb N}
\newcommand{\GAS}{\mathrm{GAS}}
\newcommand{\HLS}{\mathrm{HLS}}
\newcommand{\HGAS}{\mathrm{HGAS}}
\newcommand{\RE}{\mathrm{RE}}
\newcommand{\dd}{\,\mathrm d}
\newcommand{\Lean}{\textnormal{\textsc{Lean}}}
\newcommand{\Mathlib}{\textnormal{\textsc{Mathlib}}}
\newcommand{\leanname}[1]{\path{#1}}
\newcommand{\leanmodule}[1]{\path{#1}}
\newcommand{\manyone}{\mathrel{\leq_{\mathrm m}}}
\newcommand{\norm}[1]{\left\lVert#1\right\rVert}

\title{\textbf{Lyapunov Stability under Cubic Veronese Lifting}\\
\large A mathematical guide to the Lean formalization of shadowing,
effective cubicization, and fixed-resource transfer}
\author{Formalization project companion}
\date{8 October 2026}

\begin{document}
\maketitle

\begin{abstract}
This document explains the thirteen \Lean{} modules that extend the original
formalization of fixed-resource homogeneous global asymptotic stability.
The extension proves the Lyapunov-stability half of the cubic Veronese
reduction.  For each fixed source dimension \(n\ge2\) and each fixed
\(k\ge1\), one primitive-recursive map sends raw rational codes for
degree-\(2k+1\) homogeneous fields to raw rational codes for cubic fields in
dimension \(\binom{n+k-1}{k}\).  The same map preserves and reflects both
global asymptotic stability and Lyapunov stability with forward existence.
The development also proves the semantic shadowing theorem, degree and
dimension padding, and transfer of Lyapunov-stability hardness to every
larger fixed dimension and every odd degree at least three.

The formal proof follows the mathematical argument through a bounded
angular cocycle, exponential attraction to the Veronese cone, and
finite-horizon backward shadowing.  It does not formalize Korda's
undecidability constructions; their conclusions enter the final transfer
theorems as explicit hypotheses.  It also does not formalize a compiler
uniform in the varying pair \((n,k)\).  The orbit-shadowing declaration uses
the coordinate sup norm, rather than stating verbatim the manuscript's
weighted-\(Q\)-norm estimate; the two norms are equivalent in the fixed
finite-dimensional space, and the formal estimate is the one used to prove
the radial-cocycle bound.  These distinctions are part of the theorem
statements below, not hidden implementation details.
\end{abstract}

\tableofcontents

\section{Purpose, scope, and headline result}
\label{sec:scope}

The original companion
\leanmodule{docs/hgas_sigma1_complete_formalization_companion.tex}
documents the fixed-pair and cofinal \(\Sigma^0_1\)-completeness proof for
homogeneous GAS.  The present document is a second historical artifact.  It
covers the later Lyapunov-stability and shadowing extension without rewriting
the earlier GAS account.  The extension adds thirteen modules; the preceding
824 modules remain unchanged.

The main formal statement concerns two languages of raw rational polynomial
field codes.  The predicate \(\HGAS_{\Q}(n,d)\) requires a valid homogeneous
degree-\(d\) code whose realized field has a globally asymptotically stable
origin.  The predicate \(\HLS_{\Q}(n,d)\) requires the same syntactic validity
and the trajectory-model counterpart of Korda's Lyapunov-stability notion:
all sufficiently small initial conditions admit a solution for every forward
time, and every such solution stays small.  For polynomial fields, standard
local uniqueness identifies this notion mathematically with the usual
maximal-solution formulation.  That bridge is not packaged as a separate
Lean theorem.  Invalid homogeneity codes are negative instances of both
languages.

Put
\[
  m(n,k)=\#\{\alpha\in\N^n:|\alpha|=k\}
        =\binom{n+k-1}{k}.
\]
The \Lean{} definition is
\leanname{HGAS.Veronese.cubicLiftCodeDimension n k}.

\begin{theorem}[The formalized fixed-parameter cubicization]
\label{thm:headline}
Fix \(n,k\in\N\) with \(n\ge2\) and \(k\ge1\).  There is a computable map
\[
  \mathcal C_{n,k}:
  \operatorname{Code}(n)\longrightarrow\operatorname{Code}(m(n,k))
\]
such that every raw code \(F\) satisfies
\begin{align*}
 F\in\HGAS_{\Q}(n,2k+1)
   &\Longleftrightarrow
 \mathcal C_{n,k}(F)\in\HGAS_{\Q}(m(n,k),3),\\
 F\in\HLS_{\Q}(n,2k+1)
   &\Longleftrightarrow
 \mathcal C_{n,k}(F)\in\HLS_{\Q}(m(n,k),3).
\end{align*}
The map \(\mathcal C_{n,k}\) is the same in the two equivalences.
\end{theorem}

The exact endpoint is
\leanname{HGAS.Propagation.exists_computable_cubicization} in
\leanmodule{HGAS/Propagation/CofinalLyapunovReduction.lean}.  The explicit
map is
\leanname{HGAS.Veronese.effectiveCubicLSReductionCode}; its stronger
implementation theorem proves primitive recursiveness, not merely
computability.

The hypothesis \(n\ge2\) belongs to the total raw-code theorem.  The semantic
equivalence between a valid source field and its stabilized cubic lift has no
such hypothesis.  The restriction appears because the malformed-input branch
uses the two-coordinate unstable cubic field
\[
   \dot y_0=y_0y_1^2,\qquad \dot y_i=0\quad(i\ne0).
\]
Section~\ref{sec:codes} explains this point precisely.

\begin{corollary}[Cofinal fixed-resource transfer]
\label{cor:cofinal}
Fix \(n_0\ge2\) and \(k_0\ge1\), and put \(m_0=m(n_0,k_0)\).  If
\(\HLS_{\Q}(n_0,2k_0+1)\) is many-one hard, then
\(\HLS_{\Q}(N,d)\) is many-one hard for every \(N\ge m_0\) and every odd
\(d\ge3\).  Each target language is therefore undecidable and its complement
is not recursively enumerable.  If only undecidability of the source is
assumed, undecidability still transfers to every such target.
\end{corollary}

This corollary is unconditional as a transfer principle.  Its application to
Korda's source family is conditional in \Lean{} because the external
undecidability reduction of \cite{Korda2026} is not formalized.

\section{The semantic language and the homogeneous scaling lemma}
\label{sec:semantics}

\subsection{Forward trajectories and Lyapunov stability}

The basic dynamical object is a curve \(\gamma:\R\to E\) which starts at
\(x_0\) and solves \(\dot\gamma=F(\gamma)\) on \([0,\infty)\).  It is represented
by \leanname{HGAS.IsForwardTrajectory F gamma x0}.  Defining the curve on all
of \(\R\) is a convenient ambient representation; only its restriction to
nonnegative time is constrained.

\begin{definition}[Lyapunov stability with forward existence]
Let \(E\) be a real normed space.  The origin is Lyapunov stable with forward
existence for \(F:E\to E\) if, for every \(\varepsilon>0\), there is
\(\delta>0\) such that, whenever \(\norm{x_0}<\delta\),
\begin{enumerate}
\item at least one forward trajectory through \(x_0\) exists; and
\item every forward trajectory through \(x_0\) satisfies
      \(\norm{\gamma(t)}<\varepsilon\) for every \(t\ge0\).
\end{enumerate}
\end{definition}

This is the literal definition
\leanname{HGAS.LyapunovStableWithExistence}.  Quantifying over every forward
trajectory prevents stability from depending on a choice of solution.  For
the polynomial fields used here, local uniqueness also identifies it
mathematically with the usual maximal-solution formulation; the development
does not introduce a separate maximal-solution object.  The explicit
existence clause is essential: a definition quantifying only over already
global trajectories would hold vacuously near an initial point whose maximal
trajectory blows up.

The code language is
\[
\begin{split}
  \HLS_{\Q}(n,d)(F)\quad\Longleftrightarrow\quad&
  F\text{ decodes to a degree-\(d\) homogeneous field},\\
  &\text{the realized field is Lyapunov stable with forward existence}.
\end{split}
\]
It is \leanname{HGAS.HLSQ}.  Theorems
\leanname{HGAS.hlsQ_iff_of_homogeneous} and
\leanname{HGAS.hlsQ_of_hgasQ} expose, respectively, the semantics of valid
codes and the implication \(\HGAS\Rightarrow\HLS\).

\subsection{Stability becomes a global linear bound}

Suppose that \(F\) is homogeneous of degree \(k+1\):
\[
      F(cx)=c^{k+1}F(x)\qquad(c\in\R,\ x\in E).
\]
Homogeneity rescales both the state and time.  If \(\gamma\) is a trajectory,
then the appropriate state multiple of \(\gamma(c^k t)\) is another
trajectory.  This observation turns a local stability radius into a global
bound.

\begin{theorem}[Homogeneous scaling criterion]
\label{thm:scaling}
For a degree-\(k+1\) homogeneous field \(F\), the origin is Lyapunov stable
with forward existence if and only if there is \(C>0\) such that every
\(x_0\in E\) admits a forward trajectory and every forward trajectory through
\(x_0\) satisfies
\[
             \norm{\gamma(t)}\le C\norm{x_0}\qquad(t\ge0).
\]
\end{theorem}

The predicate on the right is
\leanname{HGAS.HasUniformLinearTrajectoryBound}; the equivalence is
\leanname{HGAS.lyapunovStableWithExistence_iff_uniformLinearTrajectoryBound}.
This is the formal analogue of Korda's Lemma~1.5
\cite[Lemma~1.5]{Korda2026}.

The nontrivial direction chooses a stability radius \(\delta\) for the unit
target ball.  For \(x_0\ne0\), scale \(x_0\) by
\(c=\delta/(2\norm{x_0})\), use stability at \(cx_0\), and scale the resulting
trajectory back.  This gives \(C=2/\delta\).  The proof treats \(x_0=0\)
separately and shows, from stability at every scale, that every trajectory
from the origin remains at the origin.  The reverse implication uses
\(\delta=\varepsilon/(2C)\).

\section{The angular criterion for Lyapunov stability}
\label{sec:angular}

Let \(E\) now be finite-dimensional and Euclidean.  For \(u\) on the unit
sphere define the radial rate and angular field
\[
  a_F(u)=\langle u,F(u)\rangle,
  \qquad Y_F(u)=F(u)-a_F(u)u.
\]
The angular flow \(\phi_\tau(u)\) remains on the sphere.  Its logarithmic
radial cocycle is
\[
  A_F(\tau,u)=\int_0^\tau a_F(\phi_s(u))\,\dd s.
\]
In the implementation, globally cut-off versions of the angular field give
a total ambient flow; the cut-off equals one on the sphere, so the displayed
objects have the intended values there.

Define boundedness of the forward angular cocycle by
\[
   \exists M\in\R\ \forall \norm{u}=1\ \forall\tau\ge0,
        \quad A_F(\tau,u)\le M.
\]
This is \leanname{HGAS.HasBoundedEuclideanAngularCocycle}.

\begin{theorem}[Angular Lyapunov criterion]
\label{thm:angular}
Let \(F:E\to E\) be \(C^1\), homogeneous of degree \(k+1\), and satisfy
\(F(0)=0\).  Then the origin is Lyapunov stable with forward existence if
and only if the forward angular cocycle is bounded uniformly over the unit
sphere.
\end{theorem}

The exact declaration is
\leanname{HGAS.lyapunovStableWithExistence_iff_boundedEuclideanAngularCocycle}.
It formalizes the Euclidean content of Korda's Lemma~2.2
\cite[Lemma~2.2]{Korda2026}.

The proof records the physical-time change rather than treating it as an
informal polar-coordinate calculation.  The angular reconstruction is
\[
       \rho(\tau)=e^{A_F(\tau,u)}\phi_\tau(u),
\]
and the physical clock has positive speed \(e^{-kA_F(\tau,u)}\).  The theorem
\leanname{HGAS.reconstruction_eq_forwardTrajectory} proves that any forward
trajectory from a unit vector, sampled at this clock, equals the
reconstruction.  A uniform linear trajectory bound therefore gives an upper
bound for \(A_F\).

Conversely, if \(A_F\le M\), the clock speed is at least \(e^{-kM}>0\) on
forward angular time.  Hence the clock reaches every nonnegative physical
time.  Inverting it constructs a global trajectory from each unit vector;
the reconstruction gives the bound \(e^M\).  Homogeneous rescaling supplies
trajectories and the same linear bound from arbitrary initial states.  The
supporting declarations are
\leanname{HGAS.exists_forwardTrajectory_of_boundedCocycle} and
\leanname{HGAS.norm_forwardTrajectory_le_of_boundedCocycle}.

Two transport results make the criterion reusable.  The declaration
\leanname{HGAS.lyapunovStableWithExistence_conjugateField_iff} proves
invariance under a continuous linear equivalence.  The declarations
\leanname{HGAS.Veronese.lyapunovStableWithExistence_iff_boundedSourceAngularCocycle}
and
\leanname{HGAS.Veronese.lyapunovStableWithExistence_iff_boundedWeightedAngularCocycle}
then express Theorem~\ref{thm:angular} in the source coordinates and in the
weighted Veronese coordinates used below.

\section{The stabilized cubic Veronese lift}
\label{sec:lift}

Fix \(n\ge1\) and \(k\ge1\).  Let
\[
 A_k=\{\alpha\in\N^n:|\alpha|=k\},\qquad
 E_k=\R^{A_k},\qquad
 \nu_k(x)=(x^\alpha)_{\alpha\in A_k}.
\]
The weighted quadratic form
\[
       Q(y)=\sum_{\alpha\in A_k}\binom{k}{\alpha}y_\alpha^2
\]
satisfies \(Q(\nu_k(x))=|x|^{2k}\).  Let \(q_r\) range over the balanced
quadratic binomials defining the real Veronese cone, and put
\[
          W(y)=\sum_r q_r(y)^2.
\]
Thus \(W=0\) is the cone and \(Q=1,\ W=0\) is its normalized core.

\subsection{A global distance estimate}

The shadowing argument needs more than the set-theoretic description of the
cone: it needs a quantitative projection estimate.

\begin{proposition}[Distance to the normalized cone]
\label{prop:distance}
There is \(C>0\) such that, for every \(y\in E_k\) with \(Q(y)=1\), there is
\(p\in E_k\) with \(Q(p)=1\), \(W(p)=0\), and
\[
          \norm{y-p}_Q^2\le C W(y).
\]
\end{proposition}

This is
\leanname{HGAS.Veronese.exists_coneDistanceSq_le_energy}.  The proof is
elementary and is fully formalized.  If a pure coordinate \(y_{ke_i}\) is
bounded away from zero, an explicit affine chart reconstructs a cone point.
A one-token transfer relation controls the chart error inductively, coordinate
by coordinate.  Normalizing the chart point to \(Q=1\) costs only a fixed
factor.  Compactness of the weighted unit sphere supplies the complementary
case: either some pure coordinate is uniformly large or \(W\) itself is
uniformly large.  The latter alternative permits any fixed normalized cone
point.  The proof therefore does not assume an unformalized tubular-neighborhood
or algebraic error-bound theorem.

\subsection{The cubic lift and its shadowing data}

Let \(F:\R^n\to\R^n\) be homogeneous of degree \(2k+1\).  The monomial
factorization of \(D\nu_k(x)F(x)\) defines a cubic field \(G:E_k\to E_k\)
such that
\[
       G(\nu_k(x))=D\nu_k(x)F(x).
\]
The stabilized cubic field is
\[
       H(y)=G(y)-K\nabla_Q W(y).
\]
It agrees with \(G\) on the cone.  On the weighted unit sphere its angular
field has the form
\[
       Y_H=Y_G-K\nabla_{S_Q}W,
\]
and its radial rate is \(a_H=a_G-4KW\).

The existing GAS formalization packages the coercivity, tangency, radial,
and gain estimates in
\leanname{HGAS.Veronese.CubicLiftStabilizationBounds}.  The new structure
\leanname{HGAS.Veronese.CubicLiftShadowingBounds} extends this package with a
constant \(L\ge0\) satisfying, on \(Q=1\),
\begin{align*}
  \norm{Y_G(y)-Y_G(y')}_\infty
       &\le L\norm{y-y'}_\infty,\\
  |a_G(y)-a_G(y')|
       &\le L\norm{y-y'}_\infty,
\end{align*}
and with the stronger gain condition
\[
       2L+1\le Kc-C.
\]
Here \(c>0\) is the transverse coercivity constant and \(C\) bounds the
tangency error.  We write
\[
       \lambda=Kc-C>0.
\]

The theorem
\leanname{HGAS.Veronese.exists_cubicLiftShadowingBounds} proves that such a
package exists for every homogeneous source field.  The supporting module
\leanmodule{HGAS/Veronese/ShadowingPrelims.lean} proves four analytic facts:
\begin{enumerate}
\item the angular field and radial reward are Lipschitz on the compact
      region \(Q\le1\);
\item the tangential penalty gradient satisfies
      \(\norm{\nabla_{S_Q}W}_Q^2\le C_W^2W\) on \(Q=1\);
\item a stabilized orbit starting in \(W\le\eta\) remains in that tube and
      satisfies \(W(t)\le W(0)e^{-\lambda t}\); and
\item stabilized and unstabilized angular flows agree on \(W=0\).
\end{enumerate}

\section{Finite-horizon shadowing}
\label{sec:shadowing}

The central obstacle is that convergence of \(W(\phi_t^H(y))\) to zero does
not identify one cone orbit which follows the stabilized orbit for all
future time.  Lyapunov stability, however, needs only a uniform upper bound
on every finite cocycle segment.  The comparison orbit may therefore depend
on the terminal horizon.

\begin{theorem}[Formal orbit-shadowing statement]
\label{thm:orbit-shadow}
Let \(b\) be a shadowing package for a homogeneous degree-\(2k+1\) source
field, with \(k\ge1\).  There is \(D\ge0\) such that, for every \(y\) with
\(Q(y)=1\) and \(W(y)\le\eta\), and every horizon \(\tau\ge0\), there is a
point \(q\) with \(Q(q)=1\) and \(W(q)=0\) such that
\[
 \norm{\phi_t^H(y)-\phi_t^H(q)}_\infty
 \le D\exp(-\lambda t/2)
 \qquad(0\le t\le\tau).
\]
\end{theorem}

This is the literal norm and flow formulation of
\leanname{HGAS.Veronese.exists_finiteHorizonOrbitShadowing}.  Because the
stabilized and unstabilized flows agree on the cone, \(\phi_t^H(q)\) may
equivalently be read as the lifted cone flow \(\phi_t^G(q)\).

The proof works backward from the endpoint.
\begin{enumerate}
\item Exponential penalty decay and the penalty-gradient bound show that
      the stabilized trajectory is an approximate \(Y_G\)-trajectory whose
      forcing error is \(O(e^{-\lambda t/2})\).
\item Proposition~\ref{prop:distance} gives a cone point \(p\) within
      \(O(e^{-\lambda\tau/2})\) of the endpoint \(\phi_\tau^H(y)\).
\item Flow \(p\) backward for time \(\tau\) on the invariant cone, and call
      the result \(q\).  Then the cone trajectory from \(q\) ends at \(p\).
\item Compare the two trajectories backward on \([0,\tau]\).  The
      \(L\)-Lipschitz separation of \(Y_G\) competes with transverse decay.
      The inequality \(\lambda\ge2L+1\) leaves a strict margin and yields
      the displayed exponential estimate, uniformly in \(\tau\).
\end{enumerate}

The radial reward is \(L\)-Lipschitz in the same sup norm, so integration of
the orbit estimate gives a horizon-independent cocycle error.

\begin{corollary}[Formal radial-cocycle shadowing]
\label{cor:cocycle-shadow}
There is \(C_*\in\R\) such that, under the hypotheses of
Theorem~\ref{thm:orbit-shadow}, one may choose \(q\) with
\[
       A_H(y,\tau)\le A_H(q,\tau)+C_*.
\]
\end{corollary}

The exact declaration is
\leanname{HGAS.Veronese.exists_finiteHorizonShadowing}.  This cocycle
inequality, rather than the choice of a particular norm in the orbit
estimate, is the fact consumed by the stability proof.

\begin{remark}[Relation to the manuscript's shadowing proposition]
The manuscript states its orbit estimate in the weighted norm
\(\norm{\cdot}_Q\).  The named \Lean{} theorem states it in the coordinate sup
norm carried by \leanname{RealVeroneseSpace}.  In the fixed finite-dimensional
space the two norms are equivalent; the development proves the required norm
comparisons internally, and the sup-norm estimate implies a weighted-norm
estimate after multiplying \(D\) by a fixed constant.  No public declaration
exports that last restatement verbatim.  Thus the formalization proves an
equivalent-strength shadowing estimate and the exact cocycle conclusion used
in the paper, but not the paper's displayed \(Q\)-norm sentence word for word.
\end{remark}

\section{The semantic Lyapunov-stability equivalence}
\label{sec:semantic-equivalence}

The angular criterion reduces Lyapunov stability to uniform upper bounds on
radial cocycles.  On the normalized cone, the lifted cocycle is exactly \(k\)
times the source cocycle.  Hence boundedness for the stabilized lift implies
boundedness for the source by restriction to the positive Veronese sheet.

For the converse, consider any stabilized angular orbit.  Before it enters
the tube \(W\le\eta\), the penalty makes its radial rate nonpositive, so this
part cannot increase the cocycle.  After entry, the tube is forward invariant
and Corollary~\ref{cor:cocycle-shadow} compares the remaining finite segment
with a cone segment up to \(C_*\).  The cone segment is \(k\) times a source
segment.  A source bound \(M\) therefore gives a stabilized bound of the form
\(\max\{0,kM+C_*\}\).

\begin{theorem}[Semantic preservation and reflection of stability]
\label{thm:semantic-ls}
Let \(F\) be homogeneous of degree \(2k+1\), with \(k\ge1\), and let \(b\)
be any valid shadowing package. Then the origin is Lyapunov stable with
forward existence for \(F\) if and only if it is Lyapunov stable with forward
existence for the stabilized cubic Veronese lift determined by \(b\).
\end{theorem}

The exact theorem is
\leanname{HGAS.Veronese.sourceLS_iff_stabilizedCubicLS}.  It is a semantic
theorem about polynomial fields and has no \(n\ge2\) assumption.  It is the
Lyapunov-stability analogue of the older GAS theorem
\leanname{HGAS.Veronese.sourceGAS_iff_stabilizedCubicGAS}.

\section{From semantic lifting to a total computable compiler}
\label{sec:codes}

\subsection{A fixed-parameter effective gain}

The semantic theorem permits any sufficiently large gain.  A reduction on
finite codes must compute one without asking whether the source is stable.
For fixed \((n,k)\), expand a homogeneous source field in the finite basis of
coordinate monomial fields.  Lipschitz constants add under sums and scale by
the absolute value of a coefficient.  Consequently there is a constant
\(B_{n,k}\) such that the needed Lipschitz constant is at most
\[
       B_{n,k}\norm{F}_1,
\]
where \(\norm{F}_1\) is the sum of the absolute values of the rational
coefficients.  This is
\leanname{HGAS.Veronese.exists_cubicLiftLipschitz_uniform}.

The development combines \(B_{n,k}\) with the fixed geometric, tangency, and
radial constants.  It chooses once and for all a natural number \(M_{n,k}>0\)
large enough that
\[
       K(F)=M_{n,k}(1+\norm{F}_1)
\]
satisfies every stabilization and shadowing inequality.  The package is
\leanname{HGAS.Veronese.FixedCubicShadowGainParameters}; its existence is
\leanname{HGAS.Veronese.nonempty_fixedCubicShadowGainParameters}.

The rational code for \(K(F)\) is
\leanname{HGAS.Veronese.effectiveCubicShadowGainCode}.  For each fixed
\((n,k)\),
\leanname{HGAS.Veronese.effectiveCubicShadowGainCode_primrec} proves that it
is primitive recursive in the input code.  The theorem
\leanname{HGAS.Veronese.exists_effectiveCubicLiftShadowingBounds} connects
the decoded gain to the analytic shadowing package.

\begin{warning}[Fixed-parameter effectivity is not uniform effectivity]
The basis Lipschitz constants and then \(M_{n,k}\) are selected using
\leanname{Classical.choice}.  Once \(n\) and \(k\) are fixed, \(M_{n,k}\) is a
constant, so hard-coding it yields a genuine primitive-recursive map of
\(F\).  The development does not give an executable procedure which receives
\((n,k)\) as variable input and computes \(M_{n,k}\).  Therefore it formalizes
the cubicization separately for every fixed pair, but it does not formalize
the manuscript's uniformity proposition.
\end{warning}

\subsection{Malformed codes and the dimension hypothesis}

The total compiler
\leanname{HGAS.Veronese.effectiveCubicLSReductionCode} first tests whether
the decoded source is homogeneous of degree \(2k+1\).  A valid code is sent to
the stabilized cubic lift with gain \(K(F)\).  An invalid code is sent to
\leanname{HGAS.Veronese.unstableCubicCode}, the code for
\(\dot y_0=y_0y_1^2\), with all other coordinates constant.

This fallback is homogeneous cubic and is neither GAS nor Lyapunov stable.
Indeed, on a trajectory with \(y_1=a>0\), the first coordinate solves
\(\dot y_0=a^2y_0\) and therefore grows exponentially.  The fully formal proof
is
\leanname{HGAS.Veronese.not_lyapunovStableWithExistence_unstableCubicCode}.

The fallback needs two target coordinates.  The theorem
\leanname{HGAS.Veronese.two_le_cubicLiftCodeDimension} obtains this from
\(n\ge2\) and \(k\ge1\).  This is the sole source of the \(n\ge2\) assumption
in the bundled raw-code HLS reduction.  A different one-dimensional negative
fallback would remove this packaging restriction; no such alternative is
part of the present formalization.

The code-level correctness theorems are
\begin{center}
\leanname{HGAS.Veronese.effectiveCubicLSReductionCode_hgasQ_iff},\\
\leanname{HGAS.Propagation.effectiveCubicLSReductionCode_hlsQ_iff}.
\end{center}
Together with primitive recursiveness they yield
\leanname{HGAS.Propagation.hlsQ_manyOneReducible_cubic} and the simultaneous
endpoint in Theorem~\ref{thm:headline}.

\section{Padding and transfer of hardness}
\label{sec:transfer}

Let \(H:\R^m\to\R^m\) be cubic, let \(p,j\in\N\), and write
\(z\in\R^p\).  Define
\[
 \widehat H(x,z)
   =\bigl(|x|^2+|z|^2\bigr)^j
      \bigl(H(x),-z_1^3,\ldots,-z_p^3\bigr).
\]
The field \(\widehat H\) has dimension \(m+p\) and degree \(3+2j\).

The proof that padding preserves Lyapunov stability has three formal pieces.
First,
\leanname{HGAS.lyapunovStableWithExistence_productField_iff} proves the
product criterion.  The independent cubic sinks are GAS, so adjoining them
preserves and reflects stability.  Second, after a Euclidean coordinate
identification, the radial factor equals one on the unit sphere.  The theorem
\leanname{HGAS.lyapunovStableWithExistence_iff_of_eqOn_sphere} proves that two
homogeneous \(C^1\) fields of possibly different degrees which agree on the
unit sphere have the same Lyapunov-stability answer.  Third, the executable
polynomial-code construction is connected to this semantic field.

The resulting code theorem is
\leanname{HGAS.DegreePaddingCode.degreeDimensionPadCode_hlsQ_iff_of_homogeneous}.
Composing it with cubicization gives
\leanname{HGAS.Propagation.hlsQ_manyOneReducible_cofinalFixedPair}:
\[
 \HLS_{\Q}(n,2k+1)\manyone\HLS_{\Q}(N,d)
\]
whenever \(n\ge2\), \(k\ge1\), \(N\ge m(n,k)\), and \(d\ge3\) is odd.

The formal predicate
\leanname{HGAS.ManyOneHard P} says that every recursively enumerable
predicate on \(\N\) computably many-one reduces to \(P\).  The general lemmas
\leanname{HGAS.ManyOneHard.not_computable} and
\leanname{HGAS.ManyOneHard.complement_not_re} prove that a many-one hard
predicate is undecidable and has a non-r.e. complement.  They do not assume
that \(P\) itself is recursively enumerable.  This distinction matters for
Lyapunov stability, whose membership in the r.e. sets is not established.

The endpoint
\leanname{HGAS.Propagation.hlsQ_manyOneHard_cofinal_of_source} takes a proof
that one source language is many-one hard and returns hardness,
undecidability, and non-r.e. complement for every cofinal target pair.  The
weaker endpoint
\leanname{HGAS.Propagation.hlsQ_not_computable_cofinal_of_source} transfers
undecidability alone.

Korda's construction supplies the source hypothesis in the manuscript
\cite[Sections~3--6]{Korda2026}.  The \Lean{} theorem does not contain that
construction.  In particular, the formalization does not independently prove
that any HLS language is many-one hard; it proves the complete compiler and
the exact conditional transfer from such a source.

\section{Paper-to-Lean crosswalk}
\label{sec:crosswalk}

Here, calling a result formalized means that the named declaration, together
with the declarations listed in its row, proves the paper result or the stated
equivalent formulation.  It does not mean that every sentence of the printed
proof has been translated line by line.

\begin{longtable}{>{\raggedright\arraybackslash}p{0.20\textwidth}
                  >{\raggedright\arraybackslash}p{0.30\textwidth}
                  >{\raggedright\arraybackslash}p{0.42\textwidth}}
\toprule
Paper result & Formal status & Lean declaration and module \\
\midrule
\endfirsthead
\toprule
Paper result & Formal status & Lean declaration and module \\
\midrule
\endhead
Homogeneous scaling lemma (Korda Lemma~1.5)
& Equivalent trajectory formulation, including forward existence and a bound
  on every global forward trajectory.  No separate maximal-solution bridge is
  exported.
& \leanname{HGAS.lyapunovStableWithExistence_iff_uniformLinearTrajectoryBound},
  \leanmodule{HGAS/Stability/KordaLyapunov.lean}. \\
\addlinespace
Angular criterion for Lyapunov stability (Korda Lemma~2.2)
& Equivalent Euclidean cocycle criterion in the same trajectory model;
  transported to the weighted form by a linear equivalence.
& \leanname{HGAS.lyapunovStableWithExistence_iff_boundedEuclideanAngularCocycle},
  \leanmodule{HGAS/Stability/AngularLyapunov.lean}; and
  \leanname{HGAS.Veronese.lyapunovStableWithExistence_iff_boundedWeightedAngularCocycle},
  \leanmodule{HGAS/Veronese/LyapunovAngularCriteria.lean}. \\
\addlinespace
Distance comparison for the Veronese cone
& Exact existential quadratic error bound in the weighted metric.
& \leanname{HGAS.Veronese.exists_coneDistanceSq_le_energy},
  \leanmodule{HGAS/Veronese/ConeDistance.lean}. \\
\addlinespace
Lipschitz and tube estimates used by shadowing
& Both the paper-style weighted estimates and the sup-norm estimates used by
  the final shadowing package are formalized; fixed-parameter effectivity uses
  the latter.
& \leanname{HGAS.Veronese.exists_lipschitz_cubicLift_angular},
  \leanmodule{HGAS/Veronese/ShadowingPrelims.lean}; and
  \leanname{HGAS.Veronese.exists_cubicLiftLipschitz_uniform},
  \leanmodule{HGAS/Computability/CubicLiftEffectiveShadowGain.lean}. \\
\addlinespace
Finite-horizon shadowing proposition
& Split into an orbit theorem and a radial-cocycle theorem.  The named orbit
  theorem uses the coordinate sup norm; the paper displays the weighted norm.
& \leanname{HGAS.Veronese.exists_finiteHorizonOrbitShadowing} and
  \leanname{HGAS.Veronese.exists_finiteHorizonShadowing},
  \leanmodule{HGAS/Veronese/CubicLiftLyapunovConjugacy.lean}. \\
\addlinespace
Cubicization: Lyapunov-stability equivalence
& The semantic theorem has no \(n\ge2\) side condition; the total raw-code
  compiler does.
& \leanname{HGAS.Veronese.sourceLS_iff_stabilizedCubicLS},
  \leanmodule{HGAS/Veronese/CubicLiftLyapunovConjugacy.lean}; and
  \leanname{HGAS.Propagation.effectiveCubicLSReductionCode_hlsQ_iff},
  \leanmodule{HGAS/Propagation/CofinalLyapunovReduction.lean}. \\
\addlinespace
Cubicization: one map for GAS and HLS
& Exact fixed-\((n,k)\) code theorem for \(n\ge2\); uniformity in varying
  \((n,k)\) is not formalized.
& \leanname{HGAS.Propagation.exists_computable_cubicization},
  \leanmodule{HGAS/Propagation/CofinalLyapunovReduction.lean}. \\
\addlinespace
Degree raising and dimension padding
& Exact HLS equivalence on valid cubic codes; combined with the older GAS
  padding result.
& \leanname{HGAS.DegreePaddingCode.degreeDimensionPadCode_hlsQ_iff_of_homogeneous},
  \leanmodule{HGAS/Stability/PaddingLyapunov.lean}. \\
\addlinespace
Cubic normal form, fixed-parameter parts (i) and (ii)
& Part (i) is an exact many-one reduction. Part (ii) is covered by the
  padding equivalence for homogeneous cubic codes composed with cubicization;
  bundled HLS cubicization requires \(n\ge2\).
& \leanname{HGAS.Propagation.hlsQ_manyOneReducible_cubic} and
  \leanname{HGAS.Propagation.hlsQ_manyOneReducible_cofinalFixedPair},
  \leanmodule{HGAS/Propagation/CofinalLyapunovReduction.lean}. \\
\addlinespace
Transfer of fixed-resource HLS hardness
& Exact conditional transfer theorem.  The source-hardness premise is
  explicit.
& \leanname{HGAS.Propagation.hlsQ_manyOneHard_cofinal_of_source},
  \leanmodule{HGAS/Propagation/CofinalLyapunovReduction.lean}. \\
\addlinespace
Consequence of Korda's main reduction
& Conditional only: Korda's construction is an external input, not a Lean
  theorem in this repository.
& Instantiate the preceding endpoint with a term of type
  \leanname{HGAS.ManyOneHard (HGAS.HLSQ n0 (2*k0+1))}. \\
\addlinespace
Dimension-five undecidability consequence
& Conditional undecidability transfer; the dimension-five source theorem is
  external.
& \leanname{HGAS.Propagation.hlsQ_not_computable_cofinal_of_source},
  \leanmodule{HGAS/Propagation/CofinalLyapunovReduction.lean}. \\
\bottomrule
\end{longtable}

\section{Inventory of the thirteen added modules}
\label{sec:inventory}

\begin{longtable}{>{\raggedright\arraybackslash}p{0.40\textwidth}
                  >{\raggedright\arraybackslash}p{0.53\textwidth}}
\toprule
Module & Principal contribution \\
\midrule
\endfirsthead
\toprule
Module & Principal contribution \\
\midrule
\endhead
\leanmodule{HGAS/Stability/KordaLyapunov.lean}
& Defines stability with forward existence and \(\HLS_{\Q}\); proves the
  homogeneous uniform-linear-bound criterion. \\
\leanmodule{HGAS/Stability/AngularLyapunov.lean}
& Proves the Euclidean angular-cocycle criterion, including the physical-clock
  reconstruction and existence argument. \\
\leanmodule{HGAS/Stability/AngularCongruence.lean}
& Proves that homogeneous fields agreeing on the unit sphere have identical
  forward angular cocycles and the same Lyapunov-stability answer. \\
\leanmodule{HGAS/Stability/LinearConjugacyLyapunov.lean}
& Proves invariance of Lyapunov stability with existence under continuous
  linear coordinate equivalences. \\
\leanmodule{HGAS/Stability/PaddingLyapunov.lean}
& Proves product, cubic-sink, degree-raising, dimension-padding, and executable
  code equivalences for HLS. \\
\leanmodule{HGAS/Veronese/LyapunovAngularCriteria.lean}
& Transports the angular criterion to the source sphere and the weighted
  Veronese sphere. \\
\leanmodule{HGAS/Veronese/ConeDistance.lean}
& Proves the global quadratic distance bound from the weighted sphere to the
  normalized real Veronese cone. \\
\leanmodule{HGAS/Veronese/ShadowingPrelims.lean}
& Supplies analytic Lipschitz bounds, the penalty-gradient estimate,
  exponential tube decay, and equality of the two flows on the cone. \\
\leanmodule{HGAS/Veronese/CubicLiftShadowingBounds.lean}
& Defines the stronger stabilization package carrying Lipschitz data and the
  gain inequality needed for shadowing. \\
\leanmodule{HGAS/Veronese/CubicLiftLyapunovConjugacy.lean}
& Proves existence of shadowing data, finite-horizon shadowing, and semantic
  preservation and reflection of Lyapunov stability. \\
\leanmodule{HGAS/Computability/ManyOneHard.lean}
& Defines many-one hardness and derives undecidability and non-r.e. complement
  without assuming r.e. membership of the target. \\
\leanmodule{HGAS/Computability/CubicLiftEffectiveShadowGain.lean}
& Builds the fixed-parameter primitive-recursive gain and total compiler,
  proves the unstable fallback negative, and establishes code-level GAS and
  conditional HLS correctness. \\
\leanmodule{HGAS/Propagation/CofinalLyapunovReduction.lean}
& Discharges the semantic premise, combines GAS and HLS in one compiler, adds
  padding, and proves cofinal hardness and undecidability transfer. \\
\bottomrule
\end{longtable}

\section{Exact limitations and nonclaims}
\label{sec:limitations}

The formalization proves the dependency chain needed for the fixed-parameter
results.  It deliberately makes the following boundaries visible.

\begin{enumerate}
\item \textbf{The raw-code bundled theorem assumes \(n\ge2\).}
  The semantic lift is valid for \(n\ge1\).  The extra assumption supports
  the two-coordinate negative fallback for malformed codes.

\item \textbf{Computability is fixed-parameter.}
  For each fixed \((n,k)\), the compiler is primitive recursive in \(F\).
  The formalization does not compute the hard-coded multiplier from a
  variable input \((n,k)\), so it does not establish the paper's uniformity
  proposition.

\item \textbf{The public orbit theorem uses the sup norm.}
  It gives an equivalent-strength finite-dimensional estimate and proves the
  required cocycle theorem, but it is not a verbatim export of the paper's
  displayed weighted-\(Q\)-norm inequality.

\item \textbf{The maximal-solution bridge is mathematical context, not an
  exported Lean theorem.}
  The development represents solutions as global forward trajectories and
  includes forward existence in the stability predicate.  Polynomial local
  uniqueness makes this equivalent to the customary maximal-solution wording,
  but the repository does not expose a standalone theorem between those two
  APIs.

\item \textbf{Korda's hardness and undecidability constructions are
  external.}
  The formalization proves trajectory-model counterparts of Korda's two
  analytic lemmas used by the cubicization, but not the register-machine
  reduction of \cite{Korda2026} or the dimension-five supplement
  \cite{KordaSupp2026}.  The final transfer declarations take the
  corresponding hardness or undecidability statement as a hypothesis.

\item \textbf{Many-one hardness is not called completeness.}
  No theorem here shows that \(\HLS_{\Q}(n,d)\) is recursively enumerable.
  The formal conclusion from Korda's source is therefore hardness,
  undecidability, and a non-r.e. complement, not
  \(\Sigma^0_1\)-completeness.

\item \textbf{Variable-resource normal forms are outside this extension.}
  The formalization does not cover the manuscript's variable-dimension
  statements, even-degree conversion, decision procedure for linear fields,
  or the consequence for GAS of fields \(H-\mathrm{id}\).
\end{enumerate}

These limitations do not weaken the fixed-dimension, fixed-degree GAS
completeness theorem documented in the earlier companion.  They delimit the
new HLS and shared-cubicization claims.

\section{Reproduction and trust boundary}
\label{sec:trust}

The repository pins Lean and Mathlib at version 4.30.0.  From the repository
root containing \leanmodule{lakefile.toml}, a fresh build is obtained with
\begin{verbatim}
lake update
lake build HGAS
\end{verbatim}
The public import root is \leanmodule{HGAS.lean}.

The new endpoints are ordinary theorems, not project axioms.  Kernel axiom
inspection of the audited endpoints reports only
\leanname{propext}, \leanname{Classical.choice}, and
\leanname{Quot.sound}.  Classical choice is used, among other places, to
select fixed analytic constants and the multiplier \(M_{n,k}\); as explained
in Section~\ref{sec:codes}, this is compatible with a computable map for each
fixed pair but does not yield a uniform executable function of the pair.

Lean's kernel checks the elaborated proof terms.  Mathlib supplies the
underlying finite-dimensional analysis, ODE, integration, polynomial, and
computability libraries \cite{Lean4,Mathlib}.  The nonroutine external
mathematical input needed to instantiate the final HLS hardness transfer is
Korda's unformalized source reduction.  Identifying its maximal-solution
wording with \Lean{}'s global-forward-trajectory predicate also uses the
standard polynomial existence-and-uniqueness context described in
Section~\ref{sec:semantics}; that API bridge is not itself exported as a
Lean theorem.

\section{Conclusion}

The extension closes the formal gap between a cubic lift that preserves GAS
and one that also preserves Lyapunov stability.  Its main new idea is the
finite-horizon comparison selected from the endpoint of each orbit segment.
This construction converts exponential transverse attraction into a uniform
radial-cocycle bound without requiring one cone orbit to shadow forever.  The
formal proof then connects the semantic argument to a primitive-recursive
compiler on raw rational codes and to cofinal fixed-resource transfer.

The resulting release should therefore present two complementary artifacts:
the earlier GAS completeness companion, preserved as a record of the original
824-module development, and this companion for the thirteen-module HLS and
shared-cubicization extension.  A coverage ledger can refer readers to both
documents while distinguishing exact statements, equivalent formulations,
conditional applications, and results not formalized.

\begin{thebibliography}{9}

\bibitem{Korda2026}
M.~Korda,
\emph{Lyapunov stability of polynomial vector fields is undecidable},
arXiv:2609.22058v2, 2026.
\url{https://arxiv.org/abs/2609.22058}.

\bibitem{KordaSupp2026}
M.~Korda,
\emph{Supplementary material for "Lyapunov stability of polynomial vector
fields is undecidable": dimension five, global asymptotic stability, and the
planar case}, 2026.
\url{https://homepages.laas.fr/mkorda/arnold_supplement.pdf}.

\bibitem{Lean4}
L.~de Moura and S.~Ullrich,
\emph{The Lean 4 theorem prover and programming language},
in \emph{Automated Deduction---CADE 28}, Lecture Notes in Computer Science
12699, Springer, 2021, pp.~625--635.
\url{https://doi.org/10.1007/978-3-030-79876-5_37}.

\bibitem{Mathlib}
The Mathlib Community,
\emph{The Lean mathematical library},
in \emph{Proceedings of the 9th ACM SIGPLAN International Conference on
Certified Programs and Proofs}, 2020, pp.~367--381.
\url{https://doi.org/10.1145/3372885.3373824}.

\end{thebibliography}

\end{document}
